\section{Das Magnetfeld}
\begin{flushleft}

\hspace{3mm} - Es gibt zwei entgegengesetzte magnetische Pole (Nord-/ Südpol) \\
\hspace{3mm} - Gleiche Pole stossen sich ab, entgegengesetzte ziehen sich an \\
\hspace{3mm} - Nord-/ Südpole tretten nie isoliert auf \\
\hspace{3mm} - Magnetische Feldlinien zeigen Nordpol $\rightarrow$ Südpol \\
\hspace{3mm} - Sie bilden immer geschlossene Schleifen \\


\subsection{Die Magnetische Kraft}
\textbf{Lorentzkraft:}
\begin{formulabox}
    \[
        \Vec{F} = q \cdot \Vec{v} \times \Vec{B}
    \]
\end{formulabox}
$\Vec{F} \perp \Vec{v} \ \& \ \Vec{B}$
\[
    [B] = Tesla = T = \frac{N/C}{m/s} = \frac{N}{A \cdot m}
\]
\hspace{5mm} Alternativ: $1Gauss = 1G = 10^{-4}T$ \\[1ex]
\hspace{5mm} $1G \approx$ Stärke des Magnetfeldes an der Erdoberfläche

\subsection{Kraft auf stromdurchflossenen Leiter}
(Ann: $\Vec{B}$ is konstant über l, keine Winkeländerung)
\[
    \Vec{F} = \left( q \cdot \Vec{v_d} \times \Vec{B} \right) \left( \frac{n}{V} \right) \cdot A \cdot l
\]
\hspace{5mm} Mit: $I = \left( \frac{n}{V} \right) \cdot q \cdot \Vec{v_d} \cdot A$ folgt: \\[1ex]
\textbf{magnetische Kraft auf Leiterabschnitt}
\begin{formulabox}
    \[
        \Vec{F} = I \cdot \Vec{l} \times \Vec{B} \hspace{5mm} \Vec{l}: \text{Längenvektor in Richtung } \Vec{j}
    \]  
\end{formulabox}
\textbf{Allgemein:} (beliebige $\Vec{B}-Felder$, infinitesimaler Leiterabschnitt) 
\begin{formulabox}
    \[
        d\Vec{F} = I \cdot d\Vec{l} \times \Vec{B}
    \]
\end{formulabox}

\subsection{Bewegung einer Punktladung \texorpdfstring{$\Vec{B}$}{B}}
Da $\Vec{F} \perp \Vec{v}$ ändert die magnetische Kraft die Richtung, nicht aber den Betrag von $\Vec{v}$ \\
\hspace{3mm} $\Rightarrow$ magnetische Kräfte, verrichten keine Arbeit an Teilchen \\[2ex]

\textbf{Spezialfall:} $\Vec{B}$ ist homogen und $\Vec{v} \perp \Vec{B}$ \\
$\Rightarrow$ Das Teilchen beschreibt eine Kreisbahn! \\[1ex]
Bewegung auf Kreisbahn:
\[
    \underbrace{q \cdot v \cdot B}_{\text{magnetische Kraft}} = \underbrace{\frac{mv^2}{r}}_{\text{ Zentripetalkraft}}
\]

\textbf{Radius des Kreises:}
\begin{formulabox}
    \[
        r = \frac{mv}{|q \cdot B|}
    \]
\end{formulabox}

\textbf{Umlaufperiode:} Zyklotronperiode
\begin{formulabox}
    \[
        T = \frac{2 \pi r}{v} = \frac{2 \pi m}{|1 \cdot B|}
    \]
\end{formulabox}

\textbf{Frequenz der Kreisbewegung:} Zyklotronfrequenz
\begin{formulabox}
    \[
        \nu = \frac{1}{T} = \frac{|q \cdot B|}{2 \pi m} \Rightarrow \omega = 2\pi\nu = \frac{|q \cdot B|}{m}
    \]
\end{formulabox}

\subsection{Das Massenspektrometer}
\textbf{Kinetische Energie} nach Beschleunigung:
\[
    \frac{1}{2} m v^2 = q U
\]
\textbf{Kreisradius:}
\[
    r = \frac{m v}{|q B|} \hspace{2mm} \Rightarrow \hspace{2mm} v^2 = \frac{r^2 q^2 B^2}{m^2}
\]
\textbf{Verhältnis $\frac{m}{q}$:}
\begin{formulabox}
    \[
        \frac{1}{2}\frac{r^2 q^2 B^2}{m} = q U \hspace{2mm} \Rightarrow \hspace{2mm} \frac{m}{q} = \frac{B^2 r^2}{2 U}
    \]
\end{formulabox}

\subsection{Das auf Leiterschleifen und Magnete ausgelöste Drehmoment}
\[
    |\Vec{F_1}| = |\Vec{F_2}| = I \cdot |\Vec{a} \times \Vec{B}| = I |\Vec{a}||\Vec{B}|
\]
\\[1ex]
\textbf{Drehmoment im Punkt P:}
\begin{formulabox}
    \[
        |\Vec{M}| = |\Vec{F_2}||\Vec{b}|\cdot \sin{\theta} = I\cdot \underbrace{|\Vec{a}||\Vec{b}|}_{|\Vec{A}|} \cdot |\Vec{B}| \cdot \sin{\theta}
    \]
\end{formulabox}

\textbf{Schleife mit mehreren Windungen $n$:}
\begin{formulabox}
    \[
        |\Vec{M}| = n \cdot I \cdot |\Vec{A}||\Vec{B}| \sin{\theta}
    \]
\end{formulabox}
    
\textbf{Definition magnetisches Moment $\Vec{\mu}$:}
\[
    \Vec{\mu} = n \cdot I \cdot \Vec{A} \hspace{5mm} [\Vec{\mu}] = Am^2
\]
\textbf{Drehmoment einer Leiterschleife:}
\begin{formulabox}
    \[
        \Vec{M} = \Vec{\mu} \times \Vec{B}
    \]
\end{formulabox}

\end{flushleft}